%% ma: L10-B14-0026
%% lop: 10
%% chuong: 5
%% bai: 14
%% dang: 10.14.4
%% loai: DS
%% muc: [NB, TH, TH, TH]
%% dap_an: "ĐĐSS"
%% nguon: "chuyen-de-cac-so-dac-trung-cua-mau-so-lieu-khong-ghep-nhom-toan-10-chuong-trinh-moi.pdf (GV. Phan Nhật Linh) (P04), Bài 03 Dạng 1 Phần II Câu 3"
%% kien_thuc: "Lớp 10 Bài 13 (số trung bình, tứ phân vị), Bài 14 (khoảng biến thiên, khoảng tứ phân vị, giá trị bất thường)"
%% trang_thai: chinh-thuc
%% ghi_chu: "Tự duyệt theo yêu cầu giáo viên 22/10/2026."
\begin{ex}
Trong một tuần, nhiệt độ cao nhất trong ngày (đơn vị: $^\circ$C) tại hai thành phố Hà Nội và Điện Biên như sau:
\begin{center}
\begin{tabular}{|c|c|c|c|c|c|c|c|}\hline
Hà Nội & $23$ & $25$ & $28$ & $28$ & $32$ & $33$ & $35$\\\hline
Điện Biên & $16$ & $24$ & $26$ & $26$ & $26$ & $27$ & $28$\\\hline
\end{tabular}
\end{center}
Xét tính đúng sai của các khẳng định sau:
\choiceTF{\True Khoảng biến thiên nhiệt độ cao nhất trong ngày của Hà Nội và Điện Biên là như nhau}{\True Nếu bỏ đi giá trị $16$ thì khoảng biến thiên của mẫu số liệu Điện Biên chỉ bằng $4$}{Khoảng tứ phân vị của mẫu số liệu Hà Nội là $\Delta_Q=3$}{Khoảng tứ phân vị của mẫu số liệu Điện Biên là $\Delta_Q=8$}
\loigiai{a) Đúng: $R_{HN}=35-23=12$; $R_{ĐB}=28-16=12$.

b) Đúng: $R=28-24=4$.

c) Sai: Hà Nội $Q_1=25$, $Q_3=33$, $\Delta_Q=8$.

d) Sai: Điện Biên $Q_1=24$, $Q_3=27$, $\Delta_Q=3$.}
\end{ex}
